\documentclass[11pt]{article}
\usepackage[a4paper,margin=1.9cm,top=2.4cm,bottom=2.2cm]{geometry}
\usepackage{amsmath,amssymb,amsfonts,fontspec,microtype,fancyhdr,lastpage,enumitem,needspace}
\usepackage[most]{tcolorbox}
\usepackage{xcolor}
\definecolor{ink}{HTML}{1F3A5F}\definecolor{soft}{HTML}{EEF3F9}\definecolor{ans}{HTML}{F3F8F1}\definecolor{ansb}{HTML}{3C7A3C}
\pagestyle{fancy}\fancyhf{}
\lhead{\small\color{ink}\textbf{Linear Algebra I} \textperiodcentered\ Week 3 --- Span, Linear Independence \& Setting Up Linear Systems}
\rhead{}\cfoot{\small Mr Malek Thiab \textperiodcentered\ Worksheet \textperiodcentered\ Page \thepage\ of \pageref{LastPage}}
\renewcommand{\headrulewidth}{0.4pt}
\setlength{\parindent}{0pt}\setlength{\parskip}{4pt}
\newcommand{\lv}[1]{\textsc{\small #1}}
\begin{document}
{\Large\bfseries\color{ink} Linear Algebra I --- Week 3}\\[2pt]
{\large\bfseries\color{ink} Span, Linear Independence \& Setting Up Linear Systems}\\[3pt]
{\small Name: \underline{\hspace{5cm}}\hfill Date: \underline{\hspace{3cm}}\qquad All work \textbf{by hand} --- no calculator, no software.}
\vspace{4pt}\hrule\vspace{8pt}

\begin{tcolorbox}[colback=soft,colframe=ink,title={\bfseries Key Facts}]
\textbf{Linear combination and span.} A linear combination of $v_1,\dots,v_k\in V$ is $a_1v_1+\dots+a_kv_k$ with $a_i\in\mathbb{F}$. $\operatorname{span}\{v_1,\dots,v_k\}$ is the set of all of them; $\operatorname{span}\varnothing=\{\mathbf 0\}$. It is a subspace, and the smallest subspace containing $v_1,\dots,v_k$.

\textbf{Setting up the system.} To test $w\in\operatorname{span}\{v_1,\dots,v_k\}$ in $\mathbb{F}^{n}$, solve $a_1v_1+\dots+a_kv_k=w$. Coordinatewise this is $n$ equations in $k$ unknowns, with augmented matrix whose columns are $v_1,\dots,v_k\mid w$. Solution exists $\iff w\in$ span.

\textbf{Linear independence.} $v_1,\dots,v_k$ are \emph{independent} if $a_1v_1+\dots+a_kv_k=\mathbf 0$ forces $a_1=\dots=a_k=0$; otherwise \emph{dependent}. Test: solve the homogeneous system; only $\mathbf 0$ $\Rightarrow$ independent, a nonzero solution $\Rightarrow$ dependent (and it gives a relation).

\textbf{Useful facts.} A set containing $\mathbf 0$ is dependent. One vector is independent iff it is nonzero. Two vectors are independent iff neither is a scalar multiple of the other. $v_1,\dots,v_k$ dependent $\iff$ some $v_j\in\operatorname{span}$ of the others. If $w\in\operatorname{span}\{v_1,\dots,v_k\}$, adding $w$ does not enlarge the span. Independence and span depend on the field: $\{(1,i),(i,-1)\}$ is dependent over $\mathbb{C}$ but independent over $\mathbb{R}$.

\textbf{Counting.} More homogeneous unknowns than equations forces a nonzero solution; so more than $n$ vectors in $\mathbb{F}^{n}$ are dependent.
\end{tcolorbox}

\newpage
\section*{\large\color{ink} Worked Examples}
\begin{tcolorbox}[colback=white,colframe=ink!60,title={\bfseries Example 1 \textperiodcentered\ Membership in a span}]Is $(4,7)\in\operatorname{span}\{(1,2),(3,5)\}$ in $\mathbb{R}^{2}$?
\par\smallskip
\textbf{Solution.} \begin{enumerate}[leftmargin=1.6em,itemsep=1pt,topsep=2pt,label=\arabic*.]\item We need numbers $a,b$ with $a(1,2)+b(3,5)=(4,7)$.
\item Expand: $a(1,2)=(a,2a)$ and $b(3,5)=(3b,5b)$. Add: $(a+3b,\ 2a+5b)$.
\item Match coordinates: Eq.\,1: $a+3b=4$; Eq.\,2: $2a+5b=7$.
\item From Eq.\,1: $a=4-3b$.
\item Substitute into Eq.\,2: $2(4-3b)+5b=7$, i.e.\ $8-6b+5b=7$, i.e.\ $8-b=7$, so $b=1$.
\item Then $a=4-3(1)=1$.
\item Check: $1(1,2)+1(3,5)=(1,2)+(3,5)=(4,7)$ \checkmark.
\end{enumerate}\textbf{Answer:} Yes: $(4,7)=1(1,2)+1(3,5)$.
\end{tcolorbox}
\begin{tcolorbox}[colback=white,colframe=ink!60,title={\bfseries Example 2 \textperiodcentered\ Membership (positive)}]Is $(1,2,3)\in\operatorname{span}\{(1,0,1),(0,1,1)\}$?
\par\smallskip
\textbf{Solution.} \begin{enumerate}[leftmargin=1.6em,itemsep=1pt,topsep=2pt,label=\arabic*.]\item We seek $a,b$ with $a(1,0,1)+b(0,1,1)=(1,2,3)$.
\item Expand: $a(1,0,1)=(a,0,a)$ and $b(0,1,1)=(0,b,b)$. Add: $(a,\ b,\ a+b)$.
\item Match coordinates: $a=1$, \ $b=2$, \ $a+b=3$.
\item The first two equations give the values $a=1$ and $b=2$ directly.
\item There are three equations for two unknowns, so check the third: $a+b=1+2=3$, which equals the required $3$ \checkmark.
\item All three equations hold, so the vector is in the span.
\end{enumerate}\textbf{Answer:} Yes: $(1,2,3)=1(1,0,1)+2(0,1,1)$.
\end{tcolorbox}
\begin{tcolorbox}[colback=white,colframe=ink!60,title={\bfseries Example 3 \textperiodcentered\ Membership (negative)}]Is $(1,2,4)\in\operatorname{span}\{(1,0,1),(0,1,1)\}$?
\par\smallskip
\textbf{Solution.} \begin{enumerate}[leftmargin=1.6em,itemsep=1pt,topsep=2pt,label=\arabic*.]\item As in the previous example, $a(1,0,1)+b(0,1,1)=(a,\ b,\ a+b)$.
\item Matching $(1,2,4)$ requires $a=1$, \ $b=2$, \ $a+b=4$.
\item The first two equations force $a=1$ and $b=2$.
\item Test the third: $a+b=1+2=3$, but the equation requires $4$. Since $3\ne4$, the system is inconsistent: no $a,b$ satisfy all three equations.
\item So $(1,2,4)$ is not a linear combination of the two vectors.
\end{enumerate}\textbf{Answer:} No.
\end{tcolorbox}
\begin{tcolorbox}[colback=white,colframe=ink!60,title={\bfseries Example 4 \textperiodcentered\ Dependence, finding a relation}]Test $\{(1,2,1),(2,1,0),(1,-1,-1)\}$ for independence.
\par\smallskip
\textbf{Solution.} \begin{enumerate}[leftmargin=1.6em,itemsep=1pt,topsep=2pt,label=\arabic*.]\item Suppose $a(1,2,1)+b(2,1,0)+c(1,-1,-1)=\mathbf 0$.
\item Expand coordinatewise: $a(1,2,1)=(a,2a,a)$, $b(2,1,0)=(2b,b,0)$, $c(1,-1,-1)=(c,-c,-c)$. Sum: $(a+2b+c,\ 2a+b-c,\ a-c)$.
\item Set each coordinate to $0$: Eq.\,1: $a+2b+c=0$; Eq.\,2: $2a+b-c=0$; Eq.\,3: $a-c=0$.
\item Eq.\,3 gives $c=a$.
\item Substitute $c=a$ into Eq.\,1: $a+2b+a=2a+2b=0$, so $b=-a$.
\item Substitute $b=-a$ and $c=a$ into Eq.\,2: $2a+(-a)-a=0$, true for every $a$. So Eq.\,2 adds no new restriction and $a$ is a free variable.
\item Choose $a=1$ (any nonzero value): $(a,b,c)=(1,-1,1)$, a nontrivial solution.
\item Check: $v_1-v_2+v_3=(1,2,1)-(2,1,0)+(1,-1,-1)=(1-2+1,\ 2-1-1,\ 1-0-1)=(0,0,0)$ \checkmark.
\end{enumerate}\textbf{Answer:} Linearly dependent: $v_1-v_2+v_3=\mathbf 0$.
\end{tcolorbox}
\begin{tcolorbox}[colback=white,colframe=ink!60,title={\bfseries Example 5 \textperiodcentered\ Independence in $P_{2}$}]Show $\{1+x,\ x+x^{2},\ 1+x^{2}\}$ is linearly independent in $P_{2}(\mathbb{R})$.
\par\smallskip
\textbf{Solution.} \begin{enumerate}[leftmargin=1.6em,itemsep=1pt,topsep=2pt,label=\arabic*.]\item Suppose $a(1+x)+b(x+x^{2})+c(1+x^{2})$ is the zero polynomial.
\item Expand and collect powers of $x$: $a+ax+bx+bx^{2}+c+cx^{2}=(a+c)+(a+b)x+(b+c)x^{2}$.
\item A polynomial is zero only if all coefficients are $0$. Constant term: $a+c=0$. Coefficient of $x$: $a+b=0$. Coefficient of $x^{2}$: $b+c=0$.
\item From $a+c=0$: $a=-c$. From $a+b=0$: $b=-a=c$.
\item Substitute into $b+c=0$: $c+c=2c=0$, so $c=0$.
\item Then $a=-c=0$ and $b=c=0$. Only the trivial solution exists, so the set is linearly independent.
\end{enumerate}\textbf{Answer:} Linearly independent.
\end{tcolorbox}
\begin{tcolorbox}[colback=white,colframe=ink!60,title={\bfseries Example 6 \textperiodcentered\ The field matters}]Is $\{(1,i),(i,-1)\}$ linearly independent in $\mathbb{C}^{2}$ over $\mathbb{C}$? Over $\mathbb{R}$?
\par\smallskip
\textbf{Solution.} \begin{enumerate}[leftmargin=1.6em,itemsep=1pt,topsep=2pt,label=\arabic*.]\item Write $v_1=(1,i)$ and $v_2=(i,-1)$. First look for a complex multiple: $i\cdot v_1=(i\cdot1,\ i\cdot i)=(i,\ i^{2})=(i,-1)=v_2$ (using $i^{2}=-1$).
\item Over $\mathbb{C}$: $i\,v_1-1\cdot v_2=\mathbf 0$ is a relation with coefficients $i$ and $-1$, not both zero. So the set is dependent over $\mathbb{C}$.
\item Over $\mathbb{R}$: suppose $a(1,i)+b(i,-1)=(0,0)$ with real $a,b$. Expand: $(a+bi,\ ai-b)=(0,0)$.
\item First coordinate $a+bi=0$: a complex number is $0$ only when both real and imaginary parts are $0$, so $a=0$ and $b=0$.
\item Only the trivial solution exists over $\mathbb{R}$, so the set is independent over $\mathbb{R}$.
\end{enumerate}\textbf{Answer:} Dependent over $\mathbb{C}$; independent over $\mathbb{R}$.
\end{tcolorbox}
\begin{tcolorbox}[colback=white,colframe=ink!60,title={\bfseries Example 7 \textperiodcentered\ Proof: dependence means one vector is redundant}]Prove: if $v_1,\dots,v_k$ are linearly dependent, then some $v_j$ is a linear combination of the others.
\par\smallskip
\textbf{Solution.} \begin{enumerate}[leftmargin=1.6em,itemsep=1pt,topsep=2pt,label=\arabic*.]\item By the definition of dependence there are scalars $a_1,\dots,a_k$, \emph{not all zero}, with $a_1v_1+\dots+a_kv_k=\mathbf 0$.
\item Since not all the $a_i$ are zero, pick an index $j$ with $a_j\neq0$.
\item Isolate the $j$-th term: move all other terms to the right side, $a_jv_j=-\sum_{i\neq j}a_iv_i$.
\item Divide both sides by $a_j$. This is legal because $a_j\neq0$ and every nonzero element of a field has an inverse: $v_j=-\sum_{i\neq j}\dfrac{a_i}{a_j}\,v_i$.
\item The right side is a linear combination of the vectors $v_i$ with $i\ne j$, so $v_j$ lies in their span.
\end{enumerate}\textbf{Answer:} $v_j\in\operatorname{span}\{v_i:i\neq j\}$.
\end{tcolorbox}
\begin{tcolorbox}[colback=white,colframe=ink!60,title={\bfseries Example 8 \textperiodcentered\ Describing a span by an equation}]Find the condition on $(a,b,c)$ for $(a,b,c)\in\operatorname{span}\{(1,1,0),(0,1,1)\}$.
\par\smallskip
\textbf{Solution.} \begin{enumerate}[leftmargin=1.6em,itemsep=1pt,topsep=2pt,label=\arabic*.]\item We need numbers $x,y$ with $x(1,1,0)+y(0,1,1)=(a,b,c)$.
\item Expand: $x(1,1,0)=(x,x,0)$ and $y(0,1,1)=(0,y,y)$. Add: $(x,\ x+y,\ y)$.
\item Match coordinates: $x=a$, \ $x+y=b$, \ $y=c$.
\item The first and third equations give $x=a$ and $y=c$.
\item Substitute into the middle equation: $a+c=b$, i.e.\ $a-b+c=0$.
\item Converse: if $a-b+c=0$ then $b=a+c$ and $a(1,1,0)+c(0,1,1)=(a,\ a+c,\ c)=(a,b,c)$. So the condition is exactly right.
\end{enumerate}\textbf{Answer:} $a-b+c=0$
\end{tcolorbox}
\newpage
\section*{\large\color{ink} Practice Problems}
\Needspace{10\baselineskip}\subsection*{\normalsize\color{ink} Part A --- Linear Combinations and Span}
\Needspace{7\baselineskip}\begin{minipage}{\linewidth}\textbf{1.} \hfill{\scriptsize\color{ink}[Foundational]}\par\nopagebreak Write $(5,-1)$ as a linear combination of $(1,1)$ and $(1,-1)$ in $\mathbb{R}^{2}$.
\par\smallskip\begingroup\small\color{ink}\textbf{Hint.} \textit{Write $(5,-1)=a(1,1)+b(1,-1)$ and compare the two coordinates. You get two equations in $a$ and $b$; add them, then subtract them.}\endgroup
\vspace{3.0cm}\end{minipage}\par\vspace{4pt}
\Needspace{7\baselineskip}\begin{minipage}{\linewidth}\textbf{2.} \hfill{\scriptsize\color{ink}[Foundational]}\par\nopagebreak Decide whether $(2,3,4)\in\operatorname{span}\{(1,1,1),(0,1,2)\}$ in $\mathbb{R}^{3}$. If so, give the coefficients.
\par\smallskip\begingroup\small\color{ink}\textbf{Hint.} \textit{Set $(2,3,4)=a(1,1,1)+b(0,1,2)$ and compare coordinates. The first coordinate gives $a$; then check that the second and third coordinates agree on one value of $b$.}\endgroup
\vspace{3.4cm}\end{minipage}\par\vspace{4pt}
\Needspace{7\baselineskip}\begin{minipage}{\linewidth}\textbf{3.} \hfill{\scriptsize\color{ink}[Foundational]}\par\nopagebreak Decide whether $(1,0,0)\in\operatorname{span}\{(1,1,0),(0,1,1)\}$ in $\mathbb{R}^{3}$.
\par\smallskip\begingroup\small\color{ink}\textbf{Hint.} \textit{Set $(1,0,0)=a(1,1,0)+b(0,1,1)$. The first coordinate fixes $a$ and the third fixes $b$; then test the second coordinate with those values.}\endgroup
\vspace{3.4cm}\end{minipage}\par\vspace{4pt}
\Needspace{7\baselineskip}\begin{minipage}{\linewidth}\textbf{4.} \hfill{\scriptsize\color{ink}[Intermediate]}\par\nopagebreak Find $k$ such that $(1,2,k)\in\operatorname{span}\{(1,1,0),(1,-1,2)\}$.
\par\smallskip\begingroup\small\color{ink}\textbf{Hint.} \textit{Write $(1,2,k)=a(1,1,0)+b(1,-1,2)$. The first two coordinates form a system in $a,b$: solve it, then the third coordinate gives $k$.}\endgroup
\vspace{3.6cm}\end{minipage}\par\vspace{4pt}
\Needspace{7\baselineskip}\begin{minipage}{\linewidth}\textbf{5.} \hfill{\scriptsize\color{ink}[Intermediate]}\par\nopagebreak In $\mathbb{C}^{2}$ over $\mathbb{C}$, write $(1+i,\;2)$ as a linear combination of $(1,i)$ and $(1,-i)$.
\par\smallskip\begingroup\small\color{ink}\textbf{Hint.} \textit{Set $(1+i,\,2)=a(1,i)+b(1,-i)$ with $a,b\in\mathbb{C}$. The coordinates give $a+b=1+i$ and $i(a-b)=2$. Divide the second by $i$, using $\dfrac1i=-i$.}\endgroup
\vspace{4.4cm}\end{minipage}\par\vspace{4pt}
\Needspace{7\baselineskip}\begin{minipage}{\linewidth}\textbf{6.} \hfill{\scriptsize\color{ink}[Foundational]}\par\nopagebreak Show that $\operatorname{span}\{(1,2),(2,4)\}\subseteq\mathbb{R}^{2}$ is a line, and give its equation.
\par\smallskip\begingroup\small\color{ink}\textbf{Hint.} \textit{One vector is a multiple of the other, so the span is the set of all multiples of a single nonzero vector. Write a general element $t(1,2)$ and find the relation between its $x$ and $y$.}\endgroup
\vspace{3.4cm}\end{minipage}\par\vspace{4pt}
\Needspace{7\baselineskip}\begin{minipage}{\linewidth}\textbf{7.} \hfill{\scriptsize\color{ink}[Intermediate]}\par\nopagebreak Describe $\operatorname{span}\{(1,0,2),(0,1,-1)\}\subseteq\mathbb{R}^{3}$ by a single linear equation in $x,y,z$.
\par\smallskip\begingroup\small\color{ink}\textbf{Hint.} \textit{A general element is $a(1,0,2)+b(0,1,-1)=(a,\,b,\,2a-b)$. Read off $x=a$ and $y=b$, then express $z$ in terms of $x$ and $y$.}\endgroup
\vspace{4.0cm}\end{minipage}\par\vspace{4pt}
\Needspace{7\baselineskip}\begin{minipage}{\linewidth}\textbf{8.} \hfill{\scriptsize\color{ink}[Challenge]}\par\nopagebreak Prove that for any vectors $v_1,\dots,v_k$ in a vector space $V$, the set $\operatorname{span}\{v_1,\dots,v_k\}$ is a subspace of $V$.
\par\smallskip\begingroup\small\color{ink}\textbf{Hint.} \textit{Use the subspace test: the set is non-empty (it contains $\mathbf 0$) and closed under addition and scalar multiplication. Add two general linear combinations and regroup the terms by $v_i$.}\endgroup
\vspace{5.0cm}\end{minipage}\par\vspace{4pt}
\Needspace{10\baselineskip}\subsection*{\normalsize\color{ink} Part B --- Linear Independence and Dependence}
\Needspace{7\baselineskip}\begin{minipage}{\linewidth}\textbf{9.} \hfill{\scriptsize\color{ink}[Foundational]}\par\nopagebreak Is $\{(1,2),(3,6)\}$ linearly independent in $\mathbb{R}^{2}$?
\par\smallskip\begingroup\small\color{ink}\textbf{Hint.} \textit{Check whether one vector is a scalar multiple of the other, or solve $a(1,2)+b(3,6)=\mathbf 0$ and look for a solution other than $a=b=0$.}\endgroup
\vspace{2.6cm}\end{minipage}\par\vspace{4pt}
\Needspace{7\baselineskip}\begin{minipage}{\linewidth}\textbf{10.} \hfill{\scriptsize\color{ink}[Foundational]}\par\nopagebreak Is $\{(1,2),(2,3)\}$ linearly independent in $\mathbb{R}^{2}$?
\par\smallskip\begingroup\small\color{ink}\textbf{Hint.} \textit{Solve $a(1,2)+b(2,3)=\mathbf 0$ coordinate by coordinate. Is $a=b=0$ the only solution?}\endgroup
\vspace{3.0cm}\end{minipage}\par\vspace{4pt}
\Needspace{7\baselineskip}\begin{minipage}{\linewidth}\textbf{11.} \hfill{\scriptsize\color{ink}[Intermediate]}\par\nopagebreak Is $\{(1,0,1),(1,1,0),(0,1,1)\}$ linearly independent in $\mathbb{R}^{3}$?
\par\smallskip\begingroup\small\color{ink}\textbf{Hint.} \textit{Solve $a(1,0,1)+b(1,1,0)+c(0,1,1)=\mathbf 0$. Three equations in $a,b,c$: add and subtract them to isolate each unknown.}\endgroup
\vspace{4.0cm}\end{minipage}\par\vspace{4pt}
\Needspace{7\baselineskip}\begin{minipage}{\linewidth}\textbf{12.} \hfill{\scriptsize\color{ink}[Intermediate]}\par\nopagebreak Is $\{(1,2,3),(4,5,6),(7,8,9)\}$ linearly independent in $\mathbb{R}^{3}$? If not, give a nontrivial relation.
\par\smallskip\begingroup\small\color{ink}\textbf{Hint.} \textit{Compare $v_2-v_1$ with $v_3-v_2$. Rearranging that observation gives a relation among $v_1,v_2,v_3$.}\endgroup
\vspace{4.0cm}\end{minipage}\par\vspace{4pt}
\Needspace{7\baselineskip}\begin{minipage}{\linewidth}\textbf{13.} \hfill{\scriptsize\color{ink}[Challenge]}\par\nopagebreak Find all real $k$ for which $\{(1,1,1),(1,2,3),(1,3,k)\}$ is linearly dependent.
\par\smallskip\begingroup\small\color{ink}\textbf{Hint.} \textit{Solve $a(1,1,1)+b(1,2,3)+c(1,3,k)=\mathbf 0$. Use the first equation to eliminate $a$, express $b$ in terms of $c$, and ask when $c$ can be nonzero.}\endgroup
\vspace{4.6cm}\end{minipage}\par\vspace{4pt}
\Needspace{7\baselineskip}\begin{minipage}{\linewidth}\textbf{14.} \hfill{\scriptsize\color{ink}[Foundational]}\par\nopagebreak In $P_{2}(\mathbb{R})$, is $\{1+x,\ 1-x,\ 2\}$ linearly independent?
\par\smallskip\begingroup\small\color{ink}\textbf{Hint.} \textit{Compare coefficients: in $a(1+x)+b(1-x)+2c=0$ the coefficient of $1$ and of $x$ must both vanish. How many equations and how many unknowns do you have?}\endgroup
\vspace{3.0cm}\end{minipage}\par\vspace{4pt}
\Needspace{7\baselineskip}\begin{minipage}{\linewidth}\textbf{15.} \hfill{\scriptsize\color{ink}[Intermediate]}\par\nopagebreak In $\mathbb{C}^{2}$, consider $\{(1,i),\ (1+i,\ i-1)\}$. (a) Is it linearly independent over $\mathbb{C}$? (b) Is it linearly independent over $\mathbb{R}$?
\par\smallskip\begingroup\small\color{ink}\textbf{Hint.} \textit{(a) Over $\mathbb{C}$ the scalar may be complex: divide the first coordinates, then test the second. (b) Over $\mathbb{R}$, write $a(1,i)+b(1+i,\,i-1)=\mathbf 0$ with real $a,b$ and separate real and imaginary parts.}\endgroup
\vspace{5.0cm}\end{minipage}\par\vspace{4pt}
\Needspace{7\baselineskip}\begin{minipage}{\linewidth}\textbf{16.} \hfill{\scriptsize\color{ink}[Challenge]}\par\nopagebreak Let $u,v,w$ be linearly independent in a real vector space. Prove that $u+v,\ v+w,\ u+w$ are linearly independent.
\par\smallskip\begingroup\small\color{ink}\textbf{Hint.} \textit{Start from $a(u+v)+b(v+w)+c(u+w)=\mathbf 0$. Collect the coefficients of $u$, $v$ and $w$, then use the independence of $u,v,w$.}\endgroup
\vspace{5.0cm}\end{minipage}\par\vspace{4pt}
\Needspace{10\baselineskip}\subsection*{\normalsize\color{ink} Part C --- Setting Up Systems, Spanning Sets and Proofs}
\Needspace{7\baselineskip}\begin{minipage}{\linewidth}\textbf{17.} \hfill{\scriptsize\color{ink}[Foundational]}\par\nopagebreak Set up and solve the system that decides whether $(1,2,3)\in\operatorname{span}\{(1,0,1),(2,1,0),(0,1,1)\}$.
\par\smallskip\begingroup\small\color{ink}\textbf{Hint.} \textit{Write $(1,2,3)=a(1,0,1)+b(2,1,0)+c(0,1,1)$ as three equations: $a+2b=1$, $b+c=2$, $a+c=3$. Substitute to eliminate one unknown at a time.}\endgroup
\vspace{4.6cm}\end{minipage}\par\vspace{4pt}
\Needspace{7\baselineskip}\begin{minipage}{\linewidth}\textbf{18.} \hfill{\scriptsize\color{ink}[Intermediate]}\par\nopagebreak Find the condition on $(a,b,c)$ for $(a,b,c)\in\operatorname{span}\{(1,2,0),(0,1,1)\}$.
\par\smallskip\begingroup\small\color{ink}\textbf{Hint.} \textit{Set $(a,b,c)=s(1,2,0)+t(0,1,1)$. The first coordinate gives $s$ and the third gives $t$; substitute both into the second coordinate.}\endgroup
\vspace{4.4cm}\end{minipage}\par\vspace{4pt}
\Needspace{7\baselineskip}\begin{minipage}{\linewidth}\textbf{19.} \hfill{\scriptsize\color{ink}[Intermediate]}\par\nopagebreak Find a finite spanning set for $W=\{(x,y,z)\in\mathbb{R}^{3}:x+y+z=0\}$.
\par\smallskip\begingroup\small\color{ink}\textbf{Hint.} \textit{Solve the condition for one variable, for example $x=-y-z$, and write a general vector with $y,z$ free. Split it into a part that contains $y$ and a part that contains $z$.}\endgroup
\vspace{4.4cm}\end{minipage}\par\vspace{4pt}
\Needspace{7\baselineskip}\begin{minipage}{\linewidth}\textbf{20.} \hfill{\scriptsize\color{ink}[Intermediate]}\par\nopagebreak (a) Show that $\{(1,1),(1,2),(2,1)\}$ spans $\mathbb{R}^{2}$. (b) Is it linearly independent?
\par\smallskip\begingroup\small\color{ink}\textbf{Hint.} \textit{(a) For a general $(x,y)$, solve $a(1,1)+b(1,2)=(x,y)$ using only the first two vectors. (b) Compare the number of vectors with the dimension of $\mathbb{R}^2$, or find a relation among the three.}\endgroup
\vspace{5.0cm}\end{minipage}\par\vspace{4pt}
\Needspace{7\baselineskip}\begin{minipage}{\linewidth}\textbf{21.} \hfill{\scriptsize\color{ink}[Intermediate]}\par\nopagebreak Show that $\{1,\ 1+x,\ 1+x+x^{2}\}$ spans $P_{2}(\mathbb{R})$.
\par\smallskip\begingroup\small\color{ink}\textbf{Hint.} \textit{Take an arbitrary $a+bx+cx^2$ and write it as $p\cdot1+q(1+x)+r(1+x+x^2)$. Match the $x^2$ coefficient first, then $x$, then the constant.}\endgroup
\vspace{4.6cm}\end{minipage}\par\vspace{4pt}
\Needspace{7\baselineskip}\begin{minipage}{\linewidth}\textbf{22.} \hfill{\scriptsize\color{ink}[Challenge]}\par\nopagebreak Show that $\left\{\begin{pmatrix}1&0\\0&1\end{pmatrix},\begin{pmatrix}0&1\\1&0\end{pmatrix},\begin{pmatrix}0&1\\-1&0\end{pmatrix},\begin{pmatrix}1&0\\0&-1\end{pmatrix}\right\}$ spans $M_{2\times2}(\mathbb{R})$ and is linearly independent.
\par\smallskip\begingroup\small\color{ink}\textbf{Hint.} \textit{Write a general matrix $\begin{pmatrix}p&q\\r&s\end{pmatrix}$ as $\alpha A_1+\beta A_2+\gamma A_3+\delta A_4$ and compare entries: four equations in four unknowns. For independence, set the matrix equal to $0$.}\endgroup
\vspace{6.0cm}\end{minipage}\par\vspace{4pt}
\Needspace{7\baselineskip}\begin{minipage}{\linewidth}\textbf{23.} \hfill{\scriptsize\color{ink}[Challenge]}\par\nopagebreak Prove: if $w\in\operatorname{span}\{v_1,\dots,v_k\}$ then $\operatorname{span}\{v_1,\dots,v_k,w\}=\operatorname{span}\{v_1,\dots,v_k\}$.
\par\smallskip\begingroup\small\color{ink}\textbf{Hint.} \textit{Prove two inclusions. $\operatorname{span}\{v_1,\dots,v_k\}\subseteq\operatorname{span}\{v_1,\dots,v_k,w\}$ is immediate; for the reverse, write $w=\sum c_iv_i$ and substitute it into a general combination.}\endgroup
\vspace{5.0cm}\end{minipage}\par\vspace{4pt}
\Needspace{7\baselineskip}\begin{minipage}{\linewidth}\textbf{24.} \hfill{\scriptsize\color{ink}[Challenge]}\par\nopagebreak Let $v_1,\dots,v_k$ be linearly independent and $v_{k+1}\notin\operatorname{span}\{v_1,\dots,v_k\}$. Prove that $v_1,\dots,v_{k+1}$ are linearly independent.
\par\smallskip\begingroup\small\color{ink}\textbf{Hint.} \textit{Suppose $a_1v_1+\dots+a_{k+1}v_{k+1}=\mathbf 0$. Ask what happens if $a_{k+1}\neq0$ (where would $v_{k+1}$ lie?), then use the independence of $v_1,\dots,v_k$.}\endgroup
\vspace{5.4cm}\end{minipage}\par\vspace{4pt}
\end{document}
